        These attributes are not case-sensitive but will be converted to
        lower case.
         
        The following six attributes are used by {\bf output()} itself.
        
        \btabularx
        \hline
        \multicolumn{4}{|l|}{\bf output mode} \\
        \hline
        direct       & 3PL        & log          & enable direct from test               \\
        differential & netlist    & log or []log & differential                          \\
        defaultout   & 3PL        & log or []log & default for unconnected outputs       \\
        noshadow     & 3PL        & log or []log & don't duplicate output flip-flop      \\
        loc          & constraint & str or []str & locations                             \\
        \hline
        \btabularxend
        
        Where an attribute value is an array the dimension must be equal to
        the number of pins and each array member will apply to the corresponding
        output pad or buffer. A single attribute value will be applied to all pins.
        
        {\bf direct} must have been set prior to any assignment statements using the
        output mode variable. The other attributes are not used until completion of
        immediate execution (excluding postprocessing).

        {\bf direct} specifies that if possible the output 3-state enable will
        be driven directly from a conditional signal, not through an execution
        thread and WHEN block. This can significantly simplify the logic and
        reduce the probability of failing to meet timing constraints locally.
        Note that when this attribute is in effect the 3-state enable is
        immediately active following configuration loading regardless of any 3PL
        startup mechanisms - see {\bf
        start()}(\autorefph{sec:ipstart}).
        
        {\bf direct} cannot be passed via iopins() as it cannot be expanded to
        an array, being a single attribute applying to the whole variable and not
        to individual outputs.
        
        {\bf differential} indicates that the output pads are differential.
        Two locations must be supplied for each data bit (see {\bf loc} above).
        
        {\bf defaultout} specifies that unconnected outputs variables, or parts of
        output variables, should be tied high or low. If this attribute is not
        supplied then the action on encountering unconnected outputs depends on the
        global directive "unass\_out" (see \autorefph{sec:directives}). The values
        must be {\bf true} to tie the outputs high, or {\bf false} to tie low.
        
        {\bf noshadow} when true inhibits the duplication of flip-flops
        whose output is connected to an output buffer. By default any
        flip-flop whose output is connected to an output buffer is placed
        in the IO pad containing the buffer. If the output of the flip-flop
        is used elsewhere then the flip-flop is duplicated in a CLB with
        all inputs in parallel with the pad flip-flop and the output of the
        duplicate used for internal logic. If {\bf noshadow} is true then
        duplication will not be done, the flip-flop will of necessity
        reside in a CLB and its output will be routed to the output buffer.
        This will result in much worse clock-to-output time. Setting this
        attribute ensures that the pad output genuinely reflects the
        flip-flop state and not that of a duplicate but at the cost of
        delayed output.
       
        The {\bf loc} attribute is mandatory. The array dimension
        must exactly match the number of bits in the variable type. If the output
        is differential the array dimension must be twice the number of
        bits, the locations being used in pairs for the +ve and the -ve inputs
        (in that order). The location array matches the associated signal
        in order LSB to MSB for the array members and the signal bits.
        in order LSB to MSB for the array members and the signal bits.
            
        Pin strings are matched left to right to an output variable
        starting at the the least-significant bit. For a compound type pins are
        matched left to right from inner to outer nested types. For a struct the
        strings are matched left to right in field order.

        FPGA package pins are usually already defined in a platform header file,
        for example -
        \begin{lstlisting}[gobble=12, language=threepl]
            static(v, "[2][2]uint:2");
            iopins(a_pins, "[12]str",
                        {
                            "A0", "B0", "A1", "B1", "A2", "C2",
                            "A3", "B3", "A4", "B4", "C3", "C4"
                        }
                );
        \end{lstlisting}
        There might be additional attributes included for these pins.

        \begin{lstlisting}[gobble=12, language=threepl]
           output(,, v, "[2][2]uint:2", loc=a_pins[2..9]);
        \end{lstlisting}

        The resulting output connections would be -\\
        v[0][0] (LSB, MSB) is connected to pins "A1", "B1" \\
        v[0][1] (LSB, MSB) is connected to pins "A2", "C2" \\
        v[1][0] (LSB, MSB) is connected to pins "A3", "B3" \\
        v[1][1] (LSB, MSB) is connected to pins "A4", "B4" \\

        The following attribute is read-only and is created when {\bf
        output()} is called and can be retrieved from the output mode
        variable if required by using inbuilt function {\bf attribute()}.
        
        \btabularx
        \hline
        \multicolumn{4}{|l|}{\bf output mode} \\
        \hline
        ids          & 3PL & []str & element IDs       \\
        \hline
        \btabularxend
        
        {\bf ids} is a string array that contains the netlist identifiers of the
        output buffer elements. This attribute is read-only and is initialised
        when the output variable is created. This identifier can be used should
        there be a necessity to apply a constraint (other than location which is
        done via the port name) to the output buffer.

        The following attributes are required for constraint generation and are
        handled at completion of immediate execution by library procedure {\bf
        generateconstraints()} \autorefph{sec:lpgenerateconstraints}. The
        attribute values are single values which will be applied to all output
        bits.
        
        \btabularx
        \hline
        \multicolumn{3}{|l|}{\bf output mode attributes} \\
        \hline
        slew         & str or []str & slew rate         \\
        drive        & int or []int & drive current     \\
        tnm          & str or []str & timing group id   \\                  
        pullup       & log or []log & pullup            \\
        pulldown     & log or []log & pulldown          \\
        iostandard   & str or []str & voltage standard  \\
        \hline
        \btabularxend
       
        The {\bf slew} attribute is the string "fast" or "slow" and selects
        one of two slew rates for Xilinx output buffers (see Xilinx documentation).
        
        The {\bf drive} attribute is an integer which specifies the output drive
        current for Xilinx output buffers (see Xilinx documentation).
        
        {\bf tnm} is the name of a timing group in which the output pads
        are to be included for timing constraint purposes.
        
        If {\bf pullup} or {\bf pulldown} are true the output will be given
        a weak pullup or pulldown. It is not allowed for both these
        attributes to be true.
        
        {\bf iostandard} is a string specifying one of the allowed
        I/O standards (see Xilinx documentation).
        
        The following attributes are tolerated but ignored. They may be included
        as attributes for the input side of a bi-directional pin.
        
        \btabularx
        \hline
        \multicolumn{4}{|l|}{\bf input mode} \\
        \hline
        ibuf\_delay\_value & constraint & int & extra delay to IBUF                      \\
        ifd\_delay\_value  & constraint & int & extra delay to IFD                       \\
        \hline
        \btabularxend

        Any attributes that are not listed above will result in a fatal error.
